\chapter{Exact C2 source and its uniqueness}
\label{component:p6e5d41ed80b3}
\noindent\textbf{Component }\path{p6e5d41ed80b3}\par
{\small\noindent Source: \path{CORE/contact/source_formulas.tex}\par}
{\small\noindent SHA256: \path{6e5d41ed80b37604ed7b057c73fddb378b4da3c993233566bf3c50b41176d477}\par}
{\small\noindent Inventory use: The six actual native C2 source formulas used in RC3-RC4; formulas only | Actual C2 source rational coefficients; The six actual native C2 source formulas used in RC3-RC4; formulas only | Actual C2 source rational coefficients\par}
\medskip
\section{Exact C2 source and its uniqueness}\label{p6e5d41ed80b3:app:c2-source}
In this appendix, $c$ denotes the actual half-angle and
$C$ the event weight previously denoted $c$. Set
\[
 z=c^2,\quad A_c=1+zr,\quad
 h=b-a^2t,\quad g=(1+z)b-(1-r+a^2zt),\quad j=1+zr-t,
\]
\begin{align*}
\Delta={}&(1+z)zb^2+2z(1+z)br+(1+z)(1+a^2z)b\\
 &+z^2(1-a^2)r^2+2z(1+a^2z)r+1+a^2z,\\
N={}&(1+z)(b+r)(b+a^2(1-t))-a^2(1+a^2z)(1-r)t.
\end{align*}
The unchanged $d=1$, $\lambda=\beta=0$ source is
\begin{align*}
n&=\frac{(b+r)A_c(A_c+a^2zt)-r\Delta}{2r\Delta},&
m&=\frac{(1+a^2z)(1-r)A_c(A_c+a^2zt)}{2g\Delta},\\
C&=\frac{A_cN}{2a^2t\Delta},&
A&=\frac{a^2A_c(b+r)(A_c+a^2zt)}{2h\Delta},\\
D&=\frac{(1+z)(b+r)(b-a^2zr)(A_c+a^2zt)}{2a^2j\Delta},&
\alpha&=\frac{z^2(b+r)N}{2(1-zh)\Delta}.
\end{align*}
Its exact contact is
\[
 (u_c,v_c;W_c,X_c,Y_c,Z_c)
   =\left(c,ac;\frac ba,-cr,c\left(\frac ba-at\right),1\right).
\]
The positive chart requires $a,c,b,t>0$, $0<r<1$, $h,g,j>0$,
$1-zh>0$, $b-a^2zr>0$, and nonnegative original event weights.
The source identities and the full independent-test radial certificates
are included in the accompanying inherited modules; the source-entry
proof above imposes no hidden shape-unit hypothesis.


The numerator has the useful positive identity
\[
N=(1+c^2)(b+r)h+a^2\{g+j+rt(1+a^2c^2)\}>0.
\]
It follows by substituting the three source-gap definitions; every term
is positive in the domain used by the contact-entry theorem.
For the determinant calculation use the aliases
\[
 R_2=R_m,\quad R_3=uW-Y,\quad R_5=R_D,\quad R_\alpha=vY-Z,
 \qquad(e_1,e_2,e_3,e_4,e_5)=(n,m,C,A,D).
\]
\par\noindent\textit{Source inclusion. }The original input is reproduced once as Component p61209c7390b0 (Chapter~\ref{component:p61209c7390b0}).\par
